\renewcommand{\theposttitle}{Why you accomplish less than you want...but more than you think}
\renewcommand{\thepostdate}{2024-02-27}
\vspace*{\fill}
\begin{center}
\begin{minipage}{0.92\textwidth}
\small
A simple trick to elevate the impact of a bar graph: \\
\textbf{Substitute the bars in the graph with a symbol. }
\end{minipage}
\end{center}
\vspace{4mm}
\begin{center}
\includegraphics[width=\textwidth, height=0.24\textheight, keepaspectratio]{C:/Users/James/Documents/GitHub/VisualIdeasBook/data/images/45ec8553203d0b9d00bb824e2d917fcb.png}
\end{center}
\vspace{2mm}
\begin{center}
\begin{minipage}{0.92\textwidth}
\small
Version 1
\end{minipage}
\end{center}
\vspace{5mm}
\begin{center}
\includegraphics[width=\textwidth, height=0.24\textheight, keepaspectratio]{C:/Users/James/Documents/GitHub/VisualIdeasBook/data/images/f8d2d5bb1e2895700c7a713743ab5558.png}
\end{center}
\vspace{2mm}
\begin{center}
\begin{minipage}{0.92\textwidth}
\small
In version 1, I clarified this idea I had about how deeply you have to understand something to teach it. It got the message across but one of my core creative practices is to \textbf{keep pushing}.

How can I make this: \\
• clearer?\\
• more insightful? \\
• more memorable?

That’s when I came up with version 2 to add another layer of nuance to the visual: substitute the bars with books.

This technique comes from my filmmaking background and I call it: \textbf{superimposition}. It’s one of the tools I teach in my \textbf{Thinking in Visual Metaphors} course.
\end{minipage}
\end{center}
\vspace*{\fill}